\section*{Capítulo \ref{ch:b2:orbital-forcing} --- El forzamiento orbital y las glaciaciones}

\begin{solution}{exo:b2:orbital-forcing:1}
(a) $1.0167/0.9833 = 1.0340$: el afelio está un \qty{3.4}{\percent} más lejos que el perihelio.
(b) $a(1 + e) - a(1 - e) = 2ea = 0.0334\,a$, un \qty{3.34}{\percent} de $a$, cerca del 3.4 por ciento
que se suele citar. (c) La luz solar disminuye como el inverso del cuadrado de la distancia: $1.0340^2 = 1.069$,
alrededor de un \qty{7}{\percent} más de luz solar en el perihelio.
\end{solution}

\begin{solution}{exo:b2:orbital-forcing:2}
(a) Hoy $0.017236 \times \sin\ang{101.37} = 0.017236 \times 0.9804 = 0.01690$; hace \num{10000} años
$0.019971 \times \sin\ang{293.86} = 0.019971 \times (-0.9145) = -0.01826$; hace \num{6000} años
$0.019249 \times \sin\ang{359.99} \approx 0$, los valores de las tablas. (b) La longitud del Sol en el
perihelio es $\varpi + \ang{180}$: \ang{281.37} hoy, \ang{473.86} $-$ \ang{360} $=$ \ang{113.86}
hace \num{10000} años, \ang{179.99} hace \num{6000} años. (c) Hace \num{10000} años el perihelio
caía \ang{24} después del solsticio de junio, el más cercano a él de los tres, en el verano boreal;
hace \num{6000} años caía en el equinoccio de septiembre, donde la \omterm{def:b2:orbital-forcing:elements}{precesión climática} es nula.
\end{solution}

\begin{solution}{exo:b2:orbital-forcing:3}
(a) Las frecuencias se suman: $1/25.77 + 1/112 = 0.03880 + 0.00893 = \qty{0.04773}{kyr^{-1}}$, un periodo
de $1/0.04773 = 21.0$ mil años. (b) El resultado cae dentro del intervalo de los periodos principales,
de 19.0 a 23.7 mil años. Un solo número no puede reproducirlos: el giro de la órbita
es a su vez la suma de varios movimientos de frecuencias distintas, cada uno de los cuales se combina con la
precesión del eje para dar un periodo de la \omterm{def:b2:orbital-forcing:elements}{precesión climática}; la media son los
\num{23000} años que se suelen citar.
\end{solution}

\begin{solution}{exo:b2:orbital-forcing:4}
$\sqrt{1 - 0.0167^2} = 0.99986$. Para \ang{22.1}: $1360 \times 0.37622/(\pi \times 0.99986) =
\qty{162.9}{W.m^{-2}}$; para \ang{24.5}: $1360 \times 0.41469/(\pi \times 0.99986) =
\qty{179.5}{W.m^{-2}}$. La razón es $\sin\ang{24.5}/\sin\ang{22.1} = 1.102$: la media anual en el
polo crece un \qty{10.2}{\percent}, unos \qty{16.7}{W.m^{-2}}, entre los dos extremos.
\end{solution}

\begin{solution}{exo:b2:orbital-forcing:5}
(a) En el solsticio de junio, el factor de distancia vale $(a/r)^2 = \bigl((1 - e\sin\varpi)/(1 - e^2)\bigr)^2$. Hoy, $e\sin\varpi
= 0.016898$ y $(0.983102/0.999703)^2 = 0.9671$, frente a $1 - 0.033796 = 0.9662$;
hace \num{116000} años $e\sin\varpi = 0.043939$ y $(0.956061/0.998065)^2 = 0.9176$, frente a $1 -
0.087878 = 0.9121$. (b) En el solsticio de diciembre $(a/r)^2 = \bigl((1 + e\sin\varpi)/(1 -
e^2)\bigr)^2$: hoy $1.0347$ frente a $1.0338$, hace \num{116000} años $1.0940$ frente a $1.0879$.
(c) Los términos despreciados son de orden $e^2$: $(1 - x)^2/(1 - e^2)^2 \approx 1 - 2x + x^2 + 2e^2$
con $x = e\sin\varpi$, aquí unos $3e^2$. Valen $3 \times 0.0172^2 = 0.0009$ hoy y $3
\times 0.0440^2 = 0.0058$ hace \num{116000} años, cerca de las diferencias halladas, $0.0009$ y
$0.0055$: la aproximación es seis veces peor cuando la excentricidad es dos veces y media
mayor.
\end{solution}

\begin{solution}{exo:b2:orbital-forcing:6}
En el solsticio de junio $\delta = \varepsilon = \ang{23.446}$. $\cos h_0 = -\tan\ang{65}\tan\ang{23.446}
= -2.14451 \times 0.43369 = -0.93006$, así que $h_0 = \ang{158.44} = \qty{2.7654}{rad}$. El día dura
$2 \times 158.44/15 = 21.13$~horas (la Tierra gira \ang{15} por hora). Con $\sin\ang{65} = 0.90631$,
$\cos\ang{65} = 0.42262$, $\sin\delta = 0.39788$, $\cos\delta = 0.91744$ y $\sin h_0 = 0.36742$,
el corchete vale $2.7654 \times 0.90631 \times 0.39788 + 0.42262 \times 0.91744 \times 0.36742 = 0.99720
+ 0.14246 = 1.13966$. A la distancia media $\bar Q = (1360/\pi) \times 1.13966 = 432.90 \times 1.13966
= \qty{493.4}{W.m^{-2}}$; con el factor de distancia del \cref{exo:b2:orbital-forcing:5}, $493.4 \times
0.9671 = \qty{477.1}{W.m^{-2}}$. Como comprobación, los mismos pasos a \ang{60} N dan \qty{475.96}{W.m^{-2}}, donde las tablas dan \num{475.95}.
\end{solution}

\begin{solution}{exo:b2:orbital-forcing:7}
(a) El equinoccio de marzo es $\lambda = 0$, el de septiembre $\lambda = \pi$. Con $\sqrt{1 - e^2}
\approx 1$, la duración es $\frac{T}{2\pi}\int_0^{\pi}\bigl(1 + 2e\cos(\lambda -
\varpi)\bigr)\,\mathrm d\lambda = \frac{T}{2\pi}\bigl(\pi + 2e[\sin(\lambda - \varpi)]_0^{\pi}\bigr)
= \frac{T}{2\pi}(\pi + 4e\sin\varpi) = \frac T2 + \frac{2T}{\pi}e\sin\varpi$. (b) El otro semestre
es más corto en la misma cantidad, así que la diferencia es $(4T/\pi)\,e\sin\varpi = (4 \times
365.2422/\pi) \times 0.016898 = 465.04 \times 0.016898 = 7.86$~días. (c) La primavera y el verano juntos
superan al otoño y al invierno en $4.5 + 3 = 7.5$~días; las dos diferencias citadas están redondeadas, y la
concordancia queda dentro de ese redondeo.
\end{solution}

\begin{solution}{exo:b2:orbital-forcing:8}
(a) Para $e = 0.0046$: $\sqrt{1 - e^2} = 0.999989$ y $340/0.999989 = \qty{340.004}{W.m^{-2}}$. Para
$e = 0.0565$: $\sqrt{1 - 0.003192} = 0.998403$ y $340/0.998403 = \qty{340.544}{W.m^{-2}}$. (b) La
diferencia es de \qty{0.54}{W.m^{-2}}, un \qty{0.16}{\percent}. (c) La insolación de mediados de julio a \ang{65} N
osciló en $497 - 388 = \qty{109}{W.m^{-2}}$, alrededor del \qty{25}{\percent} de su media y doscientas
veces el cambio de la media global anual: la órbita actúa sobre el clima trasladando luz solar entre
estaciones y latitudes, no cambiando el total.
\end{solution}

\begin{solution}{exo:b2:orbital-forcing:9}
Con $\omega = 2\pi/P$ y $\tau = \qty{10}{kyr}$: para $P = 23$, $\omega\tau = 2.732$, una amplitud
relativa $1/\sqrt{1 + 7.464} = 0.344$ y un retraso $\arctan(2.732)/\omega = 1.2199/0.2732 = 4.47$~kyr;
para $P = 41$, $\omega\tau = 1.532$, amplitud $0.546$, retraso $0.9926/0.1532 = 6.48$~kyr; para $P =
100$, $\omega\tau = 0.628$, amplitud $0.847$, retraso $0.5610/0.0628 = 8.93$~kyr. El casquete de hielo
favorece el forzamiento lento: sigue el \qty{85}{\percent} de un ciclo de \num{100000} años, pero solo un
tercio de un ciclo de precesión, porque no tiene tiempo de crecer o de menguar antes de que un forzamiento rápido
se invierta. Actúa como un filtro de paso bajo.
\end{solution}

\begin{solution}{exo:b2:orbital-forcing:10}
(a) La ecuación es lineal: si $V_1$ responde a $A_1\sin\omega_1t$ y $V_2$ responde a
$A_2\sin\omega_2t$, $V_1 + V_2$ responde a su suma. Por la \cref{prop:b2:orbital-forcing:lag}, $V_1$ y
$V_2$ son sinusoides de frecuencias $\omega_1$ y $\omega_2$ (una vez extinguido el transitorio), así que
la respuesta no contiene ninguna otra frecuencia. (b) Con $V \approx B_1\sin(\omega_1t - \phi_1) +
B_2\sin(\omega_2t - \phi_2)$,
\[
\begin{aligned}
  V^2 &= \tfrac12(B_1^2 + B_2^2) - \tfrac12B_1^2\cos(2\omega_1t - 2\phi_1) -
  \tfrac12B_2^2\cos(2\omega_2t - 2\phi_2)\\
  &\quad + B_1B_2\bigl[\cos\bigl((\omega_1 - \omega_2)t - \phi_1 + \phi_2\bigr) -
  \cos\bigl((\omega_1 + \omega_2)t - \phi_1 - \phi_2\bigr)\bigr],
\end{aligned}
\]
de modo que el término $kV^2$ fuerza la ecuación en las frecuencias diferencia y suma, y en las frecuencias
dobles. (c) $1/(1/19.0 - 1/23.7) = 1/(0.05263 - 0.04219) = 95.8$ mil años; el par 19.0
y 22.4 da 125 mil años. Las frecuencias diferencia de los términos de precesión caen en la
banda de la excentricidad, como debe ser, ya que la \omterm{def:b2:orbital-forcing:elements}{precesión climática} es $e\sin\varpi$. Un casquete de hielo no lineal
podría, por tanto, extraer un ritmo de \num{100000} años de la precesión; uno lineal no
puede. Los ciclos de \num{100000} años son una prueba de comportamiento no lineal, aunque este cálculo
por sí solo no muestra de cuál.
\end{solution}

\begin{solution}{exo:b2:orbital-forcing:11}
(a) Como en la demostración de la \cref{prop:b2:orbital-forcing:annual}, $(a/r)^2\,\mathrm dt =
T\,\mathrm d\lambda/(2\pi\sqrt{1 - e^2})$, y $\sin\delta = \sin\varepsilon\sin\lambda$ es positivo
para $0 < \lambda < \pi$. Así, $E = \frac{S_0T\sin\varepsilon}{2\pi\sqrt{1 - e^2}}\int_0^{\pi}\sin\lambda
\,\mathrm d\lambda = S_0T\sin\varepsilon/(\pi\sqrt{1 - e^2})$. (b) $T = 365.25 \times 86400 =
\qty{3.1558e7}{s}$, $S_0T/\pi = \qty{1.3661e10}{J.m^{-2}}$; para \ang{22.1}, $E = 1.3661 \times 10^{10}
\times 0.37622/0.99986 = \qty{5.14}{GJ.m^{-2}}$; para \ang{24.5}, $\qty{5.67}{GJ.m^{-2}}$,
un \qty{10}{\percent} más. (c) La precesión fija cómo se reparte en el tiempo la misma energía: un verano
en el perihelio es más corto e intenso, un verano en el afelio más largo y débil. La fusión solo empieza
por encima de \qty{0}{\celsius}, así que un verano corto e intenso puede fundir más que uno largo y fresco
que reciba el mismo total.
\end{solution}

\begin{solution}{exo:b2:orbital-forcing:12}
(a) $3.93\log_2(279.4/189.7) = 3.93\log_2 1.4729 = 3.93 \times 0.5586 = \qty{2.20}{W.m^{-2}}$. (b) La
subida más pequeña es de 194 a \qty{266}{ppm}: $3.93\log_2 1.3711 = 3.93 \times 0.4554 =
\qty{1.79}{W.m^{-2}}$; la más grande, de 188 a \qty{282}{ppm}: $3.93\log_2 1.5 = 3.93 \times 0.5850 =
\qty{2.30}{W.m^{-2}}$. (c) El forzamiento del dióxido de carbono es cuatro veces el mayor cambio de la
insolación global anual, \qty{0.54}{W.m^{-2}}, y actúa en todas partes y todo el año; la insolación
estival a \ang{65} N oscila unos \qty{109}{W.m^{-2}}, pero en una latitud y una estación. La
órbita es el desencadenante, a través de los veranos de los casquetes de hielo del norte; el dióxido de carbono, con la
\omterm{def:b2:precambrian-earth:ice-albedo}{retroalimentación hielo--albedo}, es un amplificador que extiende el cambio a todo el globo.
\end{solution}

\begin{solution}{pb:b2:orbital-forcing:1}
\textbf{1.} Entonces $0.0440 \times \sin\ang{92.7} = 0.0440 \times 0.99889 = 0.04395$; ahora $0.0172 \times
\sin\ang{101.4} = 0.0172 \times 0.98027 = 0.01686$, unas 2.6 veces menos.
\textbf{2.} Entonces $1 - 0.0440 = 0.956$ y $1 + 0.0440 = 1.044$; ahora $0.9828$ y $1.0172$.
\textbf{3.} En $\lambda = \varpi + \ang{180}$: entonces \ang{272.7}, ahora \ang{281.4}, ambas cerca del
solsticio de diciembre ($\lambda = \ang{270}$); entonces la Tierra pasaba por el perihelio menos de tres días
después de él, y por el afelio en el solsticio de junio.
\textbf{4.} Entonces $(1.044/0.956)^2 = 1.0921^2 = 1.193$, ahora $(1.0172/0.9828)^2 = 1.0350^2 = 1.071$:
un \qty{19}{\percent} más de luz solar en el perihelio entonces, un \qty{7}{\percent} ahora.
\textbf{5.} Menor, \ang{22.52} frente a \ang{23.45}: el Sol de verano estaba más bajo en las latitudes altas,
y las latitudes altas recibían menos luz solar en verano y a lo largo del año.
\textbf{6.} Entonces $340/\sqrt{1 - 0.001936} = 340/0.999032 = \qty{340.33}{W.m^{-2}}$; ahora $340/0.999852
= \qty{340.05}{W.m^{-2}}$. La diferencia, \qty{0.28}{W.m^{-2}}, menos de una parte por mil,
es insignificante.
\textbf{7.} En $\lambda = \ang{90}$, $\cos(\ang{90} - \varpi) = \sin\varpi$ en la
\cref{prop:b2:orbital-forcing:distance}.
\textbf{8.} Entonces $(1 - 0.001936)/(1 - 0.04395) = 0.998064/0.956049 = 1.04395$; ahora $(1 -
0.000296)/(1 - 0.01686) = 0.999704/0.983139 = 1.01685$.
\textbf{9.} Entonces $1/1.04395^2 = 0.91758$; ahora $1/1.01685^2 = 0.96714$.
\textbf{10.} Entonces $1 - 2 \times 0.04395 = 0.91210$; ahora $1 - 2 \times 0.01686 = 0.96628$: la
aproximación yerra en $0.0055$ entonces y en $0.0009$ ahora.
\textbf{11.} Entonces $1 - 0.91758 = \qty{8.2}{\percent}$ más débil, ahora un \qty{3.3}{\percent}.
\textbf{12.} En el solsticio $\delta = \varepsilon$: \ang{22.52} entonces, \ang{23.45} ahora.
\textbf{13.} Entonces $\cos h_0 = -2.14451 \times 0.41462 = -0.88916$, $h_0 = \ang{152.77} =
\qty{2.6663}{rad}$; ahora $\cos h_0 = -2.14451 \times 0.43378 = -0.93023$, $h_0 = \ang{158.47} =
\qty{2.7659}{rad}$.
\textbf{14.} $2h_0/\ang{15}$ por hora: entonces $20.37$~horas, ahora $21.13$~horas.
\textbf{15.} Entonces $2.6663 \times 0.90631 \times 0.38301 + 0.42262 \times 0.92375 \times 0.45759 =
0.92554 + 0.17864 = 1.10418$ (\num{1.10417} sin redondear), y $432.90 \times 1.10417 =
\qty{478.0}{W.m^{-2}}$; ahora $2.7659 \times 0.90631 \times 0.39795 + 0.42262 \times 0.91741 \times
0.36697 = 0.99757 + 0.14228 = 1.13985$ (\num{1.13982} sin redondear), y $432.90 \times 1.13982 = \qty{493.4}{W.m^{-2}}$.
\textbf{16.} Entonces $478.0 \times 0.91758 = \qty{438.6}{W.m^{-2}}$; ahora $493.4 \times 0.96714 =
\qty{477.2}{W.m^{-2}}$, \qty{38.6}{W.m^{-2}} más que entonces.
\textbf{17.} A mediados de julio la longitud del Sol es \ang{120}: $\sin\delta = \sin\varepsilon\sin\ang{120}$,
una declinación de unos \ang{20} en lugar de \ang{23}, de modo que el día es más corto y el Sol está más bajo: a
\ang{65} N la insolación media diaria es máxima cerca del solsticio, cuando la declinación es mayor.
\textbf{18.} Con la distancia actual, el cambio de oblicuidad por sí solo cuesta $(493.4 - 478.0) \times
0.96714 = \qty{14.9}{W.m^{-2}}$; con la oblicuidad antigua, la distancia por sí sola cuesta $478.0 \times
(0.96714 - 0.91758) = \qty{23.7}{W.m^{-2}}$. Ambos suman \qty{38.6}{W.m^{-2}}; la distancia
pesó unas tres quintas partes.
\textbf{19.} Tras un cambio brusco, el casquete de hielo ha alcanzado una fracción $1 - \mathrm e^{-t/\tau}$ de él:
$1 - \mathrm e^{-0.5} = 0.39$ al cabo de \qty{5}{kyr}, $1 - \mathrm e^{-1} = 0.63$ al cabo de \qty{10}{kyr},
$1 - \mathrm e^{-2} = 0.86$ al cabo de \qty{20}{kyr}.
\textbf{20.} $\omega\tau = 2\pi \times 10/23 = 2.732$: el casquete de hielo sigue una fracción $1/\sqrt{1 + 2.732^2} =
0.34$ del forzamiento, con $\arctan(2.732)/\omega = 4.5$~kyr de retraso.
\textbf{21.} Medio ciclo de precesión, unos 11 a \qty{12}{kyr}, es aproximadamente un tiempo de respuesta: antes de que
el casquete de hielo haya recorrido dos tercios del camino, la insolación se invierte y lo hace retroceder.
\textbf{22.} En el modelo lineal, el hielo crece y mengua al mismo ritmo para forzamientos del mismo
tamaño: un forzamiento creciente y uno decreciente dan respuestas simétricas. Una acumulación lenta seguida de un
colapso rápido requiere una respuesta que dependa del estado o del sentido del cambio, es decir,
una respuesta no lineal.
\textbf{23.} El dióxido de carbono subió entre $3.93\log_2(266/194) = 1.8$ y $3.93\log_2(282/188) =
\qty{2.3}{W.m^{-2}}$ del \omterm{def:b2:orbital-forcing:glacials}{periodo glaciar} al \omterm{def:b2:orbital-forcing:glacials}{interglaciar}: un forzamiento sobre todo el globo y todo el
año. Liberado cuando se debilita la circulación atlántica, transmite el cambio de los veranos
del norte a todo el planeta y, con el albedo del hielo que retrocede, lo amplifica.
\textbf{24.} Cuando terminó el último \omterm{def:b2:orbital-forcing:glacials}{interglaciar}, la insolación media diaria del solsticio de junio a
\ang{65} N era de unos \textbf{\qty{439}{W.m^{-2}}}, \qty{39}{W.m^{-2}} menos que hoy y su valor más
bajo del ciclo.
\end{solution}
